% !TEX TS-program = pdflatex
% !TEX encoding = UTF-8 Unicode

% This file is a template using the "beamer" package to create slides for a talk or presentation
% - Giving a talk on some subject.
% - The talk is between 15min and 45min long.
% - Style is ornate.

% MODIFIED by Jonathan Kew, 2008-07-06
% The header comments and encoding in this file were modified for inclusion with TeXworks.
% The content is otherwise unchanged from the original distributed with the beamer package.

\documentclass{beamer}


% Copyright 2004 by Till Tantau <tantau@users.sourceforge.net>.
%
% In principle, this file can be redistributed and/or modified under
% the terms of the GNU Public License, version 2.
%
% However, this file is supposed to be a template to be modified
% for your own needs. For this reason, if you use this file as a
% template and not specifically distribute it as part of a another
% package/program, I grant the extra permission to freely copy and
% modify this file as you see fit and even to delete this copyright
% notice. 


\mode<presentation>
{
  \usetheme{Warsaw}
  % or ...

  \setbeamercovered{transparent}
  % or whatever (possibly just delete it)
}


\usepackage[english]{babel}
% or whatever

\usepackage[utf8]{inputenc}
% or whatever

\usepackage{mathrsfs}

\usepackage{times}
\usepackage[T1]{fontenc}
% Or whatever. Note that the encoding and the font should match. If T1
% does not look nice, try deleting the line with the fontenc.

%%% MATH RELATED
\usepackage{amsmath,amsthm,amssymb}
\newtheorem{prob}{Problem}
\newtheorem{idea}{Idea}
\newtheorem{defn}{Definition}
\newtheorem{thm}{Theorem}
\newtheorem{conj}{Conjecture}
\newtheorem{ex}{Example}

\newcommand{\bbc}{\mathbb{C}}
\newcommand{\bbr}{\mathbb{R}}
\newcommand{\bbn}{\mathbb{N}}
\newcommand{\Ai}{\text{Ai}}
\newcommand{\Be}{\text{Be}}
\newcommand{\wt}[1]{\widetilde{#1}}

\title{Math GRE Prep Day 4} 
\subtitle
{} % (optional)

\author[W.R. Casper] % (optional, use only with lots of authors)
{W.R. Casper}
% - Use the \inst{?} command only if the authors have different
%   affiliation.

\institute[California State University Fullerton] % (optional, but mostly needed)
{
  Department of Mathematics\\
  California State University Fullerton}
% - Use the \inst command only if there are several affiliations.
% - Keep it simple, no one is interested in your street address.

\subject{Talks}
% This is only inserted into the PDF information catalog. Can be left
% out. 



% If you have a file called "university-logo-filename.xxx", where xxx
% is a graphic format that can be processed by latex or pdflatex,
% resp., then you can add a logo as follows:

% \pgfdeclareimage[height=0.5cm]{university-logo}{university-logo-filename}
% \logo{\pgfuseimage{university-logo}}



% Delete this, if you do not want the table of contents to pop up at
% the beginning of each subsection:
\AtBeginSubsection[]
{
  \begin{frame}<beamer>{Outline}
    \tableofcontents[currentsection,currentsubsection]
  \end{frame}
}


% If you wish to uncover everything in a step-wise fashion, uncomment
% the following command: 

%\beamerdefaultoverlayspecification{<+->}


\begin{document}

\begin{frame}
  \titlepage
\end{frame}

\begin{frame}{Outline}
  \tableofcontents
  % You might wish to add the option [pausesections]
\end{frame}

% Since this a solution template for a generic talk, very little can
% be said about how it should be structured. However, the talk length
% of between 15min and 45min and the theme suggest that you stick to
% the following rules:  

% - Exactly two or three sections (other than the summary).
% - At *most* three subsections per section.
% - Talk about 30s to 2min per frame. So there should be between about
%   15 and 30 frames, all told.

\section{Linear Algebra Problems}

\subsection{Basic problems}

\begin{frame}{Problem 1}
If the matrices
$$
\left(\begin{array}{ccc}
3 & -2 & -2\\
1 & -1 & -1\\
3 & -1 & -2
\end{array}\right)
,\quad\text{and}\quad
\left(\begin{array}{ccc}
 1 & a & 0\\
-1 & b & 1\\
 2 & c &-1
\end{array}\right)
$$
are inverses of each other, what is the value of a,b, and c?
\end{frame}

\begin{frame}{Problem 2}
In an homogeneous system of 5 linear equations with 7 unknowns, the rank of the coefficient matrix is 4. What is the maximum number of independent solution vectors?
\end{frame}

\begin{frame}{Problem 3}
If $A$ is a square matrix of order $n \geq 4$, and $a_{ij} = i + j$ represents the entry
in row $i$ and column $j$, then the rank of A is always
\end{frame}

\begin{frame}{Problem 4}
 Given that $S$ and $T$ are subspaces of a vector space, which of the following
are also subspaces?
$$S\cap T,\ \ S\cup T,\ \ 2S.$$
\end{frame}

\begin{frame}{Problem 5}
If $T$ is a linear transformation mapping vectors $(1, 0, 0)$, $(0, 1, 0)$ and $(0, 0, 1)$ to the vectors $(1, 2, 3)$, $(2, 3, 1)$ and $(1, 1, -2)$ respectively what vector is the image of the vector $(3, -2, 1)$ under $T$ ?
\end{frame}

\begin{frame}{Problem 6}
Define linear operator $S$ and $T$ on the $xy$-plane ($\mathbb R^2$) as follows:\\
$S$ rotates each vector $90$ degrees counterclockwise, and $T$ reflects each vector through the $y$-axis. If
$ST$ and $TS$ denote the compositions $S\circ T$ and $T\circ S$, respectively, and $I$ is the identity map, which of the following statements is true?
$$ST=I,\quad ST=-I,\quad TS=I,\quad TS=ST,\quad TS=-ST$$
\end{frame}

\begin{frame}{Problem 7}
Let $T:\mathbb R^5\rightarrow\mathbb R^3$ be a linear transformation whose kernel is a three dimensional subspace of $\mathbb R^5$. Then is $\{T(x): x\in\mathbb R^5\}$ just the origin, or a line through the origin, or a plane through the origin, all of $\mathbb R^3$, or something else?  Is it possible to tell?
\end{frame}

\begin{frame}{Problem 8}
Let $A$, $B$ and $C$ be real $2\times 2$ matrices, and let $0$ denote the $2\times 2$ zero
matrix. Which of the following statements is/are true?
\begin{enumerate}[(a)]
\item $A^2=0$ implies $A=0$
\item $AB=AC$ implies $B=C$
\item $A$ is invertible and $A=A^{-1}$ implies $A=I$ or $A=-I$?
\end{enumerate}
\end{frame}

\begin{frame}{Problem 9}
If $V$, $W$ are $2$-dimensional subspaces of $\mathbb R^4$, then what are the possible dimensions of $V\cap W$?
\end{frame}

\begin{frame}{Problem 10}
Suppose that $V$ is the vector space of real $2\times 3$ matrices. If $T$ is a linear transformation from $V$ onto $\mathbb R^4$ (ie. surjective!), then what is the dimension of the null space of $T$?
\end{frame}

\end{document}


